
\documentclass[11pt, a4paper]{article}
 
\usepackage[utf8]{inputenc}
\usepackage[T1]{fontenc}
\usepackage{geometry}
\usepackage{amsmath, amssymb, amsthm}
\usepackage{mathtools}
\usepackage{hyperref}
\usepackage{enumitem}
\usepackage{booktabs}
\usepackage{xcolor}
\usepackage{tcolorbox}
\usepackage{titlesec}
\usepackage{parskip}
 
\geometry{margin=2.5cm}
 
\definecolor{accentblue}{HTML}{1A3A5C}
\definecolor{lightblue}{HTML}{EAF2FB}
\definecolor{codegray}{HTML}{F5F5F5}
\definecolor{rulecolor}{HTML}{2E86C1}
 
\hypersetup{
    colorlinks=true,
    linkcolor=accentblue,
    urlcolor=rulecolor,
}
 
\titleformat{\section}{\large\bfseries\color{accentblue}}{}{0em}{}[\color{rulecolor}\titlerule]
\titleformat{\subsection}{\normalsize\bfseries\color{accentblue}}{}{0em}{}
 
\tcbuselibrary{skins, breakable}
\newtcolorbox{defbox}[1]{
    enhanced,
    colback=lightblue,
    colframe=rulecolor,
    fonttitle=\bfseries,
    title=#1,
    breakable
}
 
\newtheorem{example}{Example}
\newtheorem{definition}{Definition}
 
% ---------- document ----------
\begin{document}
 
\begin{center}
    {\LARGE \textbf{Unification in Formal Verification}}\\[0.4em]
    {\small\color{gray} Formal Verification Intro Series}
\end{center}
 
\vspace{0.5em}
\hrule height 2pt
\vspace{1.2em}
 
% ----------------------------------------------------------------
\section{What is Unification?}
 
Unification is the process of finding a \textbf{substitution} $\sigma$ that makes two
symbolic expressions \emph{syntactically identical}.
 
\begin{defbox}{Core Definition}
Given two terms $s$ and $t$ built from variables, constants, and function symbols,
a \textbf{unifier} is a substitution $\sigma$ such that
\[
    \sigma(s) \;=\; \sigma(t).
\]
If such a $\sigma$ exists, $s$ and $t$ are said to be \emph{unifiable}.
The \textbf{most general unifier} (MGU) $\hat\sigma$ is the unique (up to renaming)
unifier from which every other unifier can be obtained by further substitution.
\end{defbox}
 
\subsection{Intuitive Example}
 
Let $s = f(x, b)$ and $t = f(a, y)$, where $a, b$ are constants and $x, y$ are variables.
 
The substitution $\sigma = \{x \mapsto a,\; y \mapsto b\}$ gives
\[
    \sigma(s) = f(a,b) = \sigma(t). \qquad \checkmark
\]
This $\sigma$ is the MGU. There is no strictly more general solution because both
variables are forced.
 
\vspace{0.5em}
\noindent\textbf{Failure case.} $f(a, x)$ and $f(b, x)$ cannot be unified: no substitution
can make two distinct constants $a \neq b$ equal.
 
% ----------------------------------------------------------------
\section{Why Do We Need Unification?}
 
Unification sits at the heart of automated reasoning. It answers the fundamental question:
\emph{can these two pieces of symbolic knowledge be ``matched up'' to derive new facts?}
 
\begin{itemize}[leftmargin=*, itemsep=4pt]
    \item \textbf{Automation of logical deduction.}
          Rule application in logic requires matching a rule's precondition against a
          current goal or fact. Unification does this matching \emph{with variables},
          which is far more powerful than simple pattern matching.
 
    \item \textbf{Decidability \& completeness.}
          Robinson's 1965 result showed that first-order unification is decidable and
          that the MGU always exists when a unifier exists. This underpins the
          \emph{completeness} of resolution-based provers.
 
    \item \textbf{Efficiency.}
          The nearly-linear Martelli–Montanari (1982) and Paterson–Wegman (1978)
          algorithms make unification practical in real tools.
\end{itemize}
 
% ----------------------------------------------------------------
\section{The Unification Algorithm (Robinson / Martelli--Montanari)}
 
The algorithm works on a set of \emph{equations} (a \emph{unification problem})
$U = \{s_1 \doteq t_1, \ldots, s_n \doteq t_n\}$ and transforms it by repeatedly
applying the rules below until the set is in \emph{solved form} or failure is detected.
 
\begin{center}
\renewcommand{\arraystretch}{1.4}
\begin{tabular}{lll}
\toprule
\textbf{Rule} & \textbf{Condition} & \textbf{Action} \\
\midrule
\textsc{Delete}    & $t \doteq t$                           & Remove the equation. \\
\textsc{Decompose} & $f(\bar{s}) \doteq f(\bar{t})$         & Replace with $s_i \doteq t_i$ for each $i$. \\
\textsc{Conflict}  & $f(\bar{s}) \doteq g(\bar{t}),\ f\neq g$ & \textbf{Fail}. \\
\textsc{Orient}    & $t \doteq x$, $t$ not a variable       & Replace with $x \doteq t$. \\
\textsc{Eliminate} & $x \doteq t$, $x \notin \mathrm{vars}(t)$ & Apply $\{x \mapsto t\}$ everywhere. \\
\textsc{OccursCheck} & $x \doteq t$, $x \in \mathrm{vars}(t)$ & \textbf{Fail}. \\
\bottomrule
\end{tabular}
\end{center}
 
\vspace{0.5em}
\noindent The \textsc{OccursCheck} rule prevents cyclic terms such as $x = f(x)$.
It is crucial for soundness but is sometimes omitted for performance (as in many
Prolog implementations).
 
\subsection{Complexity}
 
\begin{itemize}[leftmargin=*, itemsep=3pt]
    \item \textbf{Naive Robinson:} $O(2^n)$ in the worst case.
    \item \textbf{Martelli--Montanari / Paterson--Wegman:} $O(n \cdot \alpha(n))$, nearly linear.
\end{itemize}
 
% ----------------------------------------------------------------
\section{Applications in Formal Verification \& Related Fields}
 
\subsection{1. Logic Programming (Prolog, Datalog)}
Resolution in Prolog is driven entirely by unification. Each predicate call attempts
to unify the call's arguments with the head of each matching clause, instantiating
variables as needed. Verification tools built on top of logic engines (e.g.,
tabling-based model checkers) inherit this dependency.
 
\subsection{2. Automated Theorem Proving (ATP)}
\textbf{Resolution provers} (E, Vampire, SPASS) use unification to identify
which clauses can be combined via the \emph{resolution rule}:
\[
    \frac{C_1 \vee A \qquad C_2 \vee \neg B}
         {(C_1 \vee C_2)\sigma}
    \quad \text{where } \sigma = \mathrm{mgu}(A, B).
\]
\textbf{Paramodulation} and \textbf{superposition} extend this with equational
reasoning—also unification-based.
 
\subsection{3. Type Inference (Hindley--Milner)}
Hindley--Milner type inference, used in ML/Haskell and in many verification tools,
reduces type checking to a unification problem over type terms. Every function
application generates a constraint; the MGU of all constraints is the inferred type.
Languages with dependent types (Agda, Lean, Coq) use \emph{higher-order unification}
for this purpose.
 
\subsection{4. SMT Solvers}
Solvers such as Z3 and CVC5 use \emph{E-unification} (unification modulo equational
theories) and \emph{matching modulo commutativity/associativity} inside their
theory combination modules (DPLL(T)). Term indexing structures (discrimination trees,
path indexing) exploit unification to retrieve candidate lemmas efficiently.
 
\subsection{5. Rewriting Systems \& Term Rewriting}
Normalisation in term-rewriting systems (used in tools like Maude and ACL2) applies
rewrite rules $l \to r$ by \emph{matching} (one-sided unification) the current term
against $l$. Completion procedures (Knuth--Bendix) additionally use unification to
detect \emph{critical pairs}.
 
\subsection{6. Model Checkers \& Symbolic Execution}
Symbolic execution tools (KLEE, JavaPathFinder) represent program states as
symbolic expressions. Checking whether two states can be merged or whether a
specification is violated reduces to unification (or more generally, constraint
satisfaction) over those expressions.
 
\subsection{7. Proof Assistants (Coq, Lean, Isabelle)}
Interactive provers rely on \emph{higher-order unification} (Huet's algorithm) for
tactics such as \texttt{apply}, \texttt{rewrite}, and \texttt{auto}. Because
higher-order unification is undecidable in general, provers restrict it to
tractable fragments (e.g., \emph{pattern unification} in Lean 4).
 
% ----------------------------------------------------------------
\section{Variants and Extensions}
 
\begin{defbox}{Key Variants}
\begin{itemize}[leftmargin=*, itemsep=4pt]
    \item \textbf{Syntactic unification} — standard, as described above.
    \item \textbf{Semantic / equational unification} — modulo a theory $E$
          (e.g., AC-unification for commutative and associative operators).
    \item \textbf{Higher-order unification} — terms include $\lambda$-abstractions;
          undecidable in general (Huet 1975), but semi-decidable procedures exist.
    \item \textbf{Disunification} — find substitutions under which terms are
          \emph{not} equal; used in constraint logic programming.
    \item \textbf{Anti-unification (generalisation)} — dual problem: find the
          \emph{least general generalisation} of two terms.
\end{itemize}
\end{defbox}
 
% ----------------------------------------------------------------
\section{Quick Reference: Unification vs.\ Matching}
 
\begin{center}
\renewcommand{\arraystretch}{1.3}
\begin{tabular}{lll}
\toprule
 & \textbf{Unification} & \textbf{Matching} \\
\midrule
Variables in & Both $s$ and $t$ & Only $s$ (pattern) \\
Substitution applied to & Both terms & Pattern only \\
Typical use & Resolution, type inference & Rewriting, Prolog head \\
Complexity & $O(n\,\alpha(n))$ & $O(n)$ \\
\bottomrule
\end{tabular}
\end{center}
 
% ----------------------------------------------------------------
\section{Summary}
 
Unification is the \emph{symbolic glue} that binds the major pillars of formal
verification together. Whether you are writing Coq proofs, building an SMT query,
or studying a resolution refutation, unification is doing the work of finding
compatible variable instantiations. Understanding the MGU, the standard algorithm,
the occurs-check, and the key variants (higher-order, equational) gives you a solid
foundation for almost every automated reasoning tool you will encounter.
 
\vspace{1em}
\begin{tcolorbox}[colback=codegray, colframe=rulecolor, title=\textbf{Further Reading}]
\begin{itemize}[leftmargin=*, itemsep=3pt, topsep=0pt]
    \item J.~A.~Robinson, \emph{A Machine-Oriented Logic Based on the Resolution Principle}, JACM 1965.
    \item A.~Martelli \& U.~Montanari, \emph{An Efficient Unification Algorithm}, ACM TOPLAS 1982.
    \item F.~Baader \& W.~Snyder, \emph{Unification Theory}, in Handbook of Automated Reasoning, 2001.
    \item G.~Huet, \emph{Unification in Typed Lambda Calculus}, in $\lambda$-Calculus and Computer Science Theory, 1975.
\end{itemize}
\end{tcolorbox}
 
\end{document}
 

